\begin{table}
    \centering
    \small
    \setlength{\tabcolsep}{2.5pt}
    \begin{tabular}{lcccccccc}
    \hline
    System & $\rho$ & $T$ & $\kappa$ & $v$ & $\mathcal{L}$ & $\mathcal{D}$ & $\hat{h}$ & $\phi/\phi_0$ \\
    & [g cm$^{-3}$] & [K] & [cm$^2$ g$^{-1}$] & [km s$^{-1}$] \\
    \hline
    PPD midplane, 0.1 au & $10^{-7}$ & $1000$ & $1$ & 94 & $3\times10^{3}$ & $0.9$ & $3.8\times10^{-6}$ & $1000$ \\
    PPD midplane, 1 au & $10^{-9}$ & $280$ & $1$ & 30 & $500$ & $0.05$ & $8.3\times10^{-6}$ & $1000$ \\
    Hot Jupiter atmosphere & $10^{-7}$ & $1500$ & $0.01$ & 3 & $1$ & $10^{-5}$ & $1.3\times10^{-5}$ & $1.0$ \\
    First hydrostatic core & $5\times10^{-12}$ & $100$ & $1$ & 1 & $380$ & $10^{-3}$ & $1.2\times10^{-4}$ & $32$ \\
    Solar core & $150$ & $1.5\times10^{7}$ & $1$ & $\approx 0$ & $10^{12}$ & -- & $4.5\times10^{-3}$ & $1.3$ \\
    PPD outer region, 100 au & $10^{-15}$ & $20$ & $1$ & 3 & $0.1$ & $10^{-6}$ & $4.8\times10^{-3}$ & $1.0$ \\
    Massive protostellar envelope & $10^{-14}$ & $100$ & $1$ & 3 & $15$ & $2\times10^{-4}$ & $0.061$ & $1.0$ \\
    PPD atmosphere, 1 au & $10^{-14}$ & $200$ & $1$ & 30 & $5\times10^{-3}$ & $5\times10^{-7}$ & $0.3$ & $1.0$ \\
    Star-forming clump & $3\times10^{-19}$ & $20$ & $1$ & 5 & $1$ & $2\times10^{-5}$ & $16$ & $1.0$ \\
    ULIRG nucleus & $2.5\times10^{-19}$ & $80$ & $1$ & 100 & $32$ & $0.01$ & $770$ & $1.0$ \\
    Massive-star envelope & $10^{-9}$ & $2\times10^{5}$ & $1$ & 10 & $100$ & $3\times10^{-3}$ & $1.2\times10^{3}$ & $1.0$ \\
    AGN dusty torus & $10^{-18}$ & $300$ & $5$ & 200 & $15$ & $0.01$ & $10^{4}$ & $1.0$ \\
    Cold molecular cloud & $2\times10^{-22}$ & $15$ & $0.1$ & 1 & $7\times10^{-5}$ & $2\times10^{-10}$ & $10^{4}$ & $1.0$ \\
    Diffuse ISM & $10^{-24}$ & $10^{4}$ & -- & 10 & $\ll 1$ & $\ll 1$ & $1.5\times10^{14}$ & $1.0$ \\
    \hline
    \end{tabular}
    \caption{Representative parameters of astrophysical systems, good to a factor of 2--3, and the gain of the ARSA, sorted by the heat-capacity ratio $\hat{h}$. Columns are, from left to right, the system being described (here PPD = protoplanetary disk, AGN = active galactic nucleus, ULIRG = ultraluminous infrared galaxy), its characteristic density $\rho$, temperature $T$, specific opacity $\kappa$, velocity $v$, optical depth across the system $\mathcal{L}$, dynamic diffusion number $\mathcal{D} = \mathcal{M}\mathcal{L}/\mathcal{C} = \mathcal{L}v/c$, ratio of radiation to gas heat capacity $\hat{h} = 4a_rT^3/(\rho c_v)$, and the reduction factor that the limiter of the ARSA sets at $\phi_0 = 10^{-3}$ in a cell of optical depth $\mathcal{L}/100$, i.e., with 100 cells across the system, for isotropic radiation in three dimensions at the RSLA step $\nu_0 = 0.9/\sqrt{3}$, $\phi/\phi_0$. We adopt $c_v = 2\times10^8$~erg~g$^{-1}$~K$^{-1}$ for ionized hydrogen and helium, $8\times10^7$ and $5\times10^7$ for warm and cold molecular gas, and $1.5\times10^8$ in the solar core.
    \label{tab:media}
    }
\end{table}
